%% notes-tex/database/ecoli-datasets/sections/caveats.tex
%% [[experiments.database.ecoli-datasets]]
%% https://github.com/Mjvolk3/torchcell/tree/main/notes-tex/database/ecoli-datasets/sections/caveats.tex
%% SOURCE: hand-authored prose; every number quoted here is a macro from
%% tables/numbers.tex, generated by
%% experiments/database/scripts/build_ecoli_genslm_benchmark_table.py

\section{What a score on these does not show\secstatus{todo}}\label{sec:caveats}

\pfstep{Gene identifiers are the released ones} Each store keys its genes as its
release does, and the \emph{Gene identifier} column of Table~\ref{tab:shortlist} is
measured from the records: b-numbers in one store, BW25113 locus tags or ECK accessions
in another, bare symbols in a third. Resolving those to the coding sequences of the
pinned assembly is a modeling-time step, by design, and its coverage must be measured
before any score is read, because a gene that fails to resolve is silently absent from
both the training and the test fold. The ECK synonym is the published one-to-one key
between the two K-12 strains and is the route for the Keio-based rows.

\pfstep{Held-out genes can leak through paralogs} A paralog of a held-out gene in the
training fold gives a sequence model a near copy of the test input. A split that
clusters genes by sequence similarity before holding out clusters removes that route; a
random split does not, and a random-split score is an upper bound.

\pfstep{The pretraining corpus is not disjoint from the test set} GenSLM was
pretrained on gene sequences from BV-BRC, a database that holds \org{E. coli} K-12
assemblies, so the genes being scored are expected, unverified for any one gene, to be
in its training data. The vector can then carry gene identity rather than a property
of the sequence, which is exactly what the one-hot control measures: an encoding that
scores no better than a one-hot vector of the same width has memorized, not
generalized.

\pfstep{Guide-level stores are excluded on purpose} The \guideStoreCount{} CRISPRi
stores that score one guide at a time are the largest \org{E. coli} rows, and the gene
each guide targets does have a coding sequence. The label, though, depends on the 20-nt
protospacer, its position in the gene and its knockdown strength, none of which the
gene's vector sees. Aggregating guides to their gene is what Wang 2018 released, and
that is the row on the list.

\pfstep{A tier is not a result} The sequence-only fit tier and the pick list say what
a sequence-only model could be given. No encoding has been scored on any of these rows
here; every expectation in Sec.~\ref{sec:benchmark} is labeled as one.
